


Guo et al. \cite{guo2025systematic}
provides a foundational threat model for the Model Context Protocol by formally categorizing the novel vulnerabilities introduced when Large Language Models integrate with third-party environments. The authors systematically analyze how the architecture's reliance on external servers creates distinct attack vectors, such as host-side prompt injection, server-side tool poisoning, and malicious data exposure. By mapping these theoretical vulnerabilities to practical exploits, the research establishes a necessary taxonomy for understanding how implicit trust between the MCP host and server can be weaponized by adversaries.

MCP-SafetyBench \cite{zong2025mcp} focuses on empirical evaluation, this research introduces a comprehensive benchmark designed specifically to test the robustness of LLMs acting as MCP hosts against malicious server environments. The authors develop a framework that simulates poisoned tools and deceptive server metadata to observe whether state-of-the-art models can distinguish between legitimate instructions and harmful commands. Through extensive testing across various tasks, the paper highlights the alarming tendency of advanced models to blindly follow malicious tool descriptions, thereby establishing a critical baseline for measuring agent compliance and safety.

MCPSecBench \cite{yang2025mcpsecbench} broadens the scope of empirical testing by introducing a large-scale security benchmark that evaluates agents across a diverse array of threat domains. Rather than focusing on a single type of exploit, the authors construct a vast dataset of test cases that measure an agent’s susceptibility to data exfiltration, unauthorized privilege escalation, and malicious code execution within the MCP ecosystem. Specifically, the benchmark categorizes these vulnerabilities into five distinct domains and deploys 182 rigorous test cases across 65 diverse simulated environments. The benchmark's results highlight a critical industry-wide gap, revealing that even the most advanced proprietary models fail to achieve complete safety, while many open-source alternatives exhibit alarmingly low defense rates. By rigorously testing multiple commercial and open-source models against this diverse suite of attacks, the research exposes systemic vulnerabilities in how current LLMs handle unverified context from external providers.


MCP Security Bench (MSB) \cite{zhang2025msb} contributes a novel, highly quantified methodology for assessing the specific risk profiles of AI agents utilizing the Model Context Protocol. The authors design the MSB framework to move beyond binary pass/fail security metrics, instead introducing granular scoring systems that evaluate exactly how and when a model's safety guardrails fail under pressure. Specifically, the benchmark systematically evaluates models against three core vulnerability categories: unauthorized resource access, sensitive data leakage, and environment disruption. Furthermore, testing across seven state-of-the-art LLMs reveals a concerning trade-off, demonstrating that models optimized heavily for usefulness often exhibit significant security degradation when exposed to compromised servers. By detailing the empirical results of these granular evaluations, the paper provides researchers with actionable insights into the specific architectural weaknesses that make current models vulnerable to context manipulation.

MCIP \cite{jing2025mcip} proposes a novel framework designed to mitigate the inherent risks of the MCP architecture. The authors introduce the Model Contextual Integrity Protocol (MCIP), a security layer that enforces strict validation and contextual boundaries between the LLM host and external servers to prevent malicious data from manipulating the agent’s reasoning process. To achieve this, the framework operates as a generalized runtime defense that filters communications and isolates user instructions from untrusted external inputs without requiring modifications to the underlying models. By demonstrating that enforcing contextual integrity significantly reduces the success rate of tool poisoning and injection attacks, the research offers a promising, actionable pathway toward securing multi-agent orchestration.